\chapter*{Notación y convenciones}
\addcontentsline{toc}{chapter}{Notación y convenciones}

\begin{table}[htbp]
\centering
\begin{tabular}{p{0.28\textwidth} p{0.64\textwidth}}
\hline
Símbolo & Significado \\
\hline
\multicolumn{2}{l}{\textit{Vectores y coordenadas}} \\
\hline
$\vb x$ & Vector de posición en $\mathbb{R}^3$. \\
$\vb r$ & Vector transversal en $\mathbb{R}^2$, con $\vb r=(x,y)$. \\
$\mathbf k$ & Vector de onda en $\mathbb{R}^3$. \\
$\mathbf k_\perp$ & Componente transversal del vector de onda, $\mathbf k_\perp=(k_x,k_y)$. \\
$k_z$ & Componente longitudinal del vector de onda. \\
\hline
\multicolumn{2}{l}{\textit{Campos eléctricos}} \\
\hline
$\mathbb E_0$ & Vector de amplitud del campo incidente (parte de polarización y módulo, sin la fase esférica). Sus componentes en la base de Fresnel son $\mathscr E_\theta$ y $\mathscr E_\phi$. \\
$\mathscr E_\theta,\;\mathscr E_\phi$ & Amplitudes escalares del campo incidente en la base de Fresnel local $(\hat p, \hat s)$. \\
$\mathbb E^{\mathscr G}$ & Campo sobre la esfera de referencia $\mathscr G$ centrada en el punto objeto $A$; satisface $\mathbb E^{\mathscr G} = \mathbb E_0/\ell_0$. \\
\texorpdfstring{$\mathbb E^{\mathscr G'}$}{E\^{}G'} & Campo sobre la esfera de referencia $\mathscr G'$ centrada en el punto imagen $A'$, tangente al dioptrio. Función pupila vectorial generalizada; satisface $\mathbb E^{\mathscr G'} = \mathbb M\,\mathbb E^{\mathscr G}$. \\
$\vb E^{\mathscr D}$ & Campo sobre la esfera $\mathscr D$ centrada en el vértice del dioptrio con radio $f$. Proporcional a la transformada de Fourier de $\mathbb E^{\mathscr G'}$. \\
$\mathbb M$ & Operador de transmisión de Fresnel; aplica $t_p$ a la componente $p$ y $t_s$ a la componente $s$. \\
$\Psi(\vb x)$ & Componente escalar del campo en el espacio tridimensional. \\
$\Psi(\vb r,z)$ & Campo escalar en coordenadas transversales y longitudinal. \\
\hline
\multicolumn{2}{l}{\textit{Coeficientes de Fresnel}} \\
\hline
$t_s,\,t_p$ & Coeficientes de transmisión de Fresnel para polarización $s$ y $p$. \\
$t_+ = t_s+t_p$ & Combinación isótropa de los coeficientes de Fresnel. \\
$t_- = t_p-t_s$ & Combinación anisotrópica; gobierna la mezcla de polarizaciones. \\
\hline
\multicolumn{2}{l}{\textit{Operadores y transformadas}} \\
\hline
$\mathscr F$ & Operador de transformada de Fourier. \\
$\mathscr F_\perp$ & Transformada de Fourier bidimensional transversal. \\
$\mathscr H_n$ & Transformada de Hankel de orden $n$. \\
$\mathscr R$ & Operador de propagación (notación general). \\
$\mathscr R_{ASM}$ & Operador de propagación del método de espectro angular. \\
\hline
\multicolumn{2}{l}{\textit{Otras cantidades}} \\
\hline
$\mathrm{NA}$ & Apertura numérica del sistema óptico. \\
$F_c$ & Razón entre la intensidad central vectorial y la escalar: $I_\text{total}(0)/I_\text{escalar}(0)$. \\
\hline
\end{tabular}
\caption{Convenciones de notación empleadas en el documento.}
\label{tab:notation_conventions}
\end{table}

% ======================
